% CPSY 1291 -- Recitation handout 12
% Compile:  latexmk -lualatex handout-12-final-review.tex
\documentclass[11pt]{article}
\usepackage{cpsy1291-handout}

\begin{document}

\handouthead{12}{Final exam review}
  {Lectures 9--19, in eight pages and sixteen problems}
  {Week 12 \ $\cdot$\ the exam is in class, Tue 12/1 --- the first day back from recess}

\section*{How to use this}
\addcontentsline{toc}{section}{How to use this}

\hl{The exam is Tuesday 12/1, the first day back from Thanksgiving recess.}
This session is held before recess because there is no session between recess and
the exam. Plan accordingly.

The exam covers Lectures 9--19. Read Section~1 first --- it is those eleven
lectures compressed into the statements you should be able to make without notes
--- then close it and do Section~2.

\begin{keybox}
\ask{The shape of the second half of this course.} Lectures 1--8 built machinery.
Lectures 9--19 spend it on one repeated question: \emph{does this model tell us
anything about the brain, and how would we know?} Almost every lecture presents
a method for answering that, then presents a reason the method is not enough.
Both halves are examinable, and the second half is where the marks are.
\end{keybox}

\section{Lectures 9--19, compressed}

\subsection{L09 --- Generalization}

Training error and test error are different quantities, and only the second one
is a measurement of the model. Three splits, three jobs; the test set is opened
once.

The \textbf{classical bargain}: more capacity lowers training error and
eventually raises test error --- the U-shaped curve. \textbf{Bias--variance}
decomposes the error into what the model class cannot express and how much its
answer moves with the sample.

Classical theory breaks in three places, and each is examinable. Networks fit
\textbf{random labels} perfectly (Zhang et al.), so capacity-based bounds cannot
explain why they generalize. \textbf{Double descent}: past the interpolation
threshold, test error falls again. \textbf{Grokking}: generalization can appear
long after training error reaches zero.

\textbf{Compositional generalization} --- systematically recombining known parts
--- is something humans do easily and networks do poorly. \textbf{Shortcut
learning}: a model solves the task using a feature that happens to correlate with
the label in your data --- texture rather than shape, a watermark rather than a
diagnosis --- and fails out of distribution. This is the recurring failure of the
second half of the course.

\subsection{L10 --- Autoencoders and sparse coding}

An autoencoder compresses and reconstructs, and is trained on its own input, so
it needs no labels. A \textbf{linear autoencoder with tied weights, trained on
centered data, recovers the PCA subspace} --- not necessarily the individual
components, but the space they span. Add nonlinearity and it can beat PCA,
because it can follow a curved manifold.

\textbf{Sparse coding} (Olshausen \& Field): reconstruct natural image patches
using as few active units as possible. The learned dictionary is
oriented, localized, bandpass --- Gabor-like, and quantitatively similar to V1
simple-cell receptive fields. This is the course's cleanest example of a
\emph{normative} explanation: a computational objective, not a mechanism,
predicting the biology.

\textbf{Efficient coding} is the reason to expect this: natural images are
redundant, and a code that removes redundancy under a cost constraint is forced
toward these filters. \textbf{Sparse autoencoders} apply the same idea to a
trained network's activations to extract more monosemantic features, which is
the current interpretability workhorse.

\subsection{L11 --- The visual system and CNNs}

The ventral stream --- retina, LGN, V1, V2, V4, IT --- is a hierarchy of
increasing receptive-field size, increasing complexity of preferred features, and
increasing tolerance to position and size. \textbf{Core object recognition} is
the behavior to explain: fast, accurate identification despite enormous image
variation.

Center-surround in retina and LGN; \textbf{orientation-tuned} simple cells in V1,
well described by Gabors; complex cells with the same orientation preference and
more position tolerance. \textbf{HMAX} stacks that: selectivity from
template-matching, invariance from max-pooling.

A CNN's convolution is template matching with \textbf{weight sharing}; pooling
gives tolerance; depth grows receptive fields. Task-optimized CNNs were the first
models to predict IT responses well --- and models can now be used to
\emph{control} neurons via synthesized images, which is a stronger test than
prediction.

But: \textbf{texture bias} (CNNs classify by texture where humans use shape),
brittle out-of-distribution behavior, and a scorecard of human signatures most
CNNs do not reproduce.

\subsection{L12 --- RSA, model--brain comparison, and decoding}

Two systems have no shared coordinates, so compare \emph{geometry}. An RDM
discards units and keeps pairwise dissimilarity; RSA is the (Spearman)
correlation between two RDMs. Interpretation requires a \textbf{noise ceiling}
and a null obtained by \textbf{permuting stimuli}.

The \textbf{mapping hypothesis}: a model layer corresponds to a brain area if a
(usually linear) map from the model's units predicts held-out responses.
Alternatives: RSA, CKA, regression-based scores, and Brain-Score as a
standardized aggregate. \textbf{The metrics disagree}, high scores can be
engineered, and at scale many models score alike --- so a single number is weak
evidence.

\textbf{Decoding} (MVPA) asks what can be read out of a population. A decoder is
a cross-validated classifier; chance depends on the class balance; decodability
does not establish that the organism uses the information. An SVM maximizes the
margin, and regularization ($C$) matters most when features outnumber examples,
which is the normal case.

\subsection{L13 --- Explainability}

Two families. \textbf{Attribution} --- saliency (gradient of a logit with respect
to the input), Grad-CAM (feature-map level) --- is fast and fragile, and can
survive randomizing the model's weights, which means it was showing the image
rather than the model. \textbf{Synthesis} --- feature visualization by optimizing
an input --- requires strong regularization to produce anything interpretable, so
the picture is partly a statement about the prior you chose.

\textbf{Maximally-activating patches} are the oldest and often most informative
method; \textbf{Network Dissection} quantifies it against labeled concepts.
Units are frequently \textbf{polysemantic}, which motivates sparse-autoencoder
features. Meaning generally lives in \emph{directions} in activation space rather
than in single neurons. Attribution can catch shortcut learning, which is its
clearest scientific use.

\subsection{L14--L15 --- Recurrence and dynamics}

An RNN carries a state: $\boldsymbol{h}_t = g(\boldsymbol{W}_{hh}\boldsymbol{h}_{t-1}
+ \boldsymbol{W}_{xh}\boldsymbol{x}_t)$, with weights shared across time.
Unrolled, it is a deep feedforward network. BPTT differentiates through the
unrolled graph, producing a \emph{product} of Jacobians --- hence vanishing and
exploding gradients. Clipping fixes exploding; \textbf{gating} (LSTM, GRU) fixes
vanishing, because the cell state passes through an addition.

Viewed as a dynamical system, the same network has \textbf{fixed points}. Stable
ones are point attractors: discrete memories and decisions. A near-unity
eigenvalue gives a \textbf{line attractor}: continuous integration and working
memory. Trained RNNs are reverse-engineered by finding fixed points numerically
and linearizing.

\textbf{Hopfield networks} store patterns via a Hebbian outer product, recall by
descending an energy function, and hold about $0.14n$ patterns before spurious
blends dominate. Energy-based models generalize this, and connect forward to
diffusion. The brain is not feedforward: recurrent circuit models (e.g.\ the
hGRU) solve tasks like Pathfinder where feedforward networks of comparable size
fail.

\subsection{L16 --- Attention and the transformer}

Self-attention: $\mathrm{softmax}(\boldsymbol{Q}\boldsymbol{K}^{\!\top}/\sqrt{d_k})\boldsymbol{V}$,
with $\boldsymbol{Q},\boldsymbol{K},\boldsymbol{V}$ all linear projections of the
same input. Every position sees every other in one step --- no long product of
Jacobians, and it parallelizes --- at the cost of $O(T^2)$ and of losing order,
which positional encodings restore. Multi-head is a reshape. A block is
attention $+$ MLP with residual connections and LayerNorm. \textbf{ViTs} treat
image patches as tokens. Attention maps show routing, not importance.

\subsection{L17 --- LLMs, scaling, and self-supervision}

Next-token prediction, trained with \textbf{teacher forcing}. Generation samples
from the predicted distribution; temperature and top-$k$ change the output
substantially with the weights fixed. \textbf{Scaling laws}: loss falls as a
power law in parameters, data and compute; Chinchilla showed models were badly
under-trained relative to their size. \textbf{In-context learning} and
\textbf{emergent abilities} are claimed; part of the emergence literature is
sensitive to the choice of metric, and the exam may ask you to say why.

\textbf{Self-supervised learning} makes the data label itself. Two families:
\textbf{contrastive} (SimCLR, InfoNCE --- augmentations of one image attract, other
images repel) and \textbf{masked prediction} (BERT, MAE --- hide part of the input
and predict it). Both learn representations that transfer, without labels, which
matters for biology because the brain has no labeled dataset.

\subsection{L18 --- Generative models}

\textbf{VAE}: an encoder outputs a distribution; the reparameterization trick
makes sampling differentiable; the ELBO is reconstruction plus a KL to the prior.
Latent directions are meaningful and steerable, and VAE latents have been related
to IT face-patch coding.

\textbf{GAN}: generator versus discriminator, a minimax game. Sharp samples,
unstable training, mode collapse. Latent directions found in a GAN can be used as
psychological instruments --- GANalyze finds a memorability direction.

\textbf{Diffusion}: add noise with a closed-form forward process, learn to
predict the noise, sample by denoising repeatedly. The training objective is a
mean squared error; the complexity is in sampling.

\subsection{L19 --- World models and frontiers}

A \textbf{world model} is an internal model of environmental dynamics that
supports prediction and mental simulation. Whether next-step prediction alone
builds a coherent one is open --- the taxi-map result suggests it often does not.
Evidence for mental simulation in primates; recurrent models learn simulation
where feedforward ones do not.

Three limits of current AI: out-of-distribution robustness and shortcut learning;
visual and abstract reasoning; data efficiency. And the question the course
closes on: \emph{does better AI mean better models of biology?} The trend is
worrying --- as benchmark performance rises, brain-predictivity has not kept
pace --- and the proposed diagnosis is to compare \emph{strategies} rather than
outputs (ClickMe, harmonization, data diets, better objectives).

\section{Problems}

\ask{1.} A network reaches 100\% training accuracy on randomly shuffled labels.
Why does this undermine capacity-based explanations of generalization?

\ask{2.} Sketch a double-descent curve. Mark the interpolation threshold and say
what is true of the model there.

\ask{3.} A classifier reaches 94\% on a held-out test set and 41\% on
photographs taken in a different hospital. Name the phenomenon and one method
from L13 that could have caught it in advance.

\ask{4.} Why does a linear autoencoder with tied weights recover the PCA
subspace, and why is ``the PCA subspace'' the right phrasing rather than ``the
principal components''?

\ask{5.} Sparse coding on natural images produces Gabor-like filters. Explain in
two sentences why this is evidence about V1 rather than a coincidence.

\ask{6.} What does max-pooling buy, what does it cost, and which cortical cell
type is it usually likened to?

\ask{7.} A model scores 0.42 by RSA against IT. Give the two numbers you need
before this means anything, and say how each is obtained.

\ask{8.} Explain why permuting the entries of an RDM's upper triangle produces a
null distribution that is too narrow.

\ask{9.} You randomize a network's weights and its saliency map is unchanged.
What follows?

\ask{10.} An RNN trains on sequences of length 20 and fails at length 300.
Name the mechanism, and say why gradient clipping does not fix it.

\ask{11.} A trained RNN has a fixed point with one Jacobian eigenvalue at 1.00
and the rest near 0.3. What computation does this support, and give a task that
would need it.

\ask{12.} Why is the $\sqrt{d_k}$ in scaled dot-product attention necessary, and
what fails without it?

\ask{13.} A colleague shows an attention map and claims it explains the model's
answer. Give two objections.

\ask{14.} Teacher forcing makes training faster and creates a problem. State the
problem.

\ask{15.} Explain posterior collapse in a VAE, and give the one-line diagnostic.

\ask{16.} State the closing question of L19 and the evidence that makes it
worrying.

\subsection*{Answers}

\ask{1.} If a network can fit arbitrary labels, its capacity is sufficient to
memorize the training set outright --- so capacity cannot be what prevents it from
memorizing when the labels are real. Whatever explains generalization must
therefore involve the data, the optimizer, or the architecture's inductive bias,
not the parameter count.

\ask{2.} Test error falls, rises to a peak at the interpolation threshold, then
falls again. At the threshold the model has just enough capacity to fit the
training data exactly --- roughly, parameters $\approx$ training examples --- and
the fit is maximally constrained and unstable.

\ask{3.} Shortcut learning: the model used a feature that correlates with the
label in the training hospital (scanner, marking, preprocessing) rather than the
pathology. Attribution --- a saliency or Grad-CAM map --- would show the model
attending to a region that could not carry the diagnosis.

\ask{4.} Minimizing squared reconstruction error through a $k$-dimensional
linear bottleneck is exactly the problem PCA solves, so the optimum spans the top
$k$ principal directions. It recovers the \emph{subspace} rather than the
components because any rotation within that subspace reconstructs equally well ---
nothing in the objective orders or aligns the axes.

\ask{5.} The filters were not fitted to neural data: they fall out of an
objective (reconstruct natural images with few active units) applied to natural
images alone, and they then match measured V1 receptive fields quantitatively.
That makes it a normative explanation --- the biology is what you would expect if
V1 were solving that problem --- rather than a curve fit.

\ask{6.} It buys tolerance to small shifts (and reduces spatial resolution
cheaply); it costs position information, which later layers cannot recover. It is
likened to the complex cell, which shares its simple cells' orientation
preference but responds across a range of positions.

\ask{7.} The \textbf{noise ceiling}, from a split-half correlation of the neural
RDMs, which bounds what any model could achieve; and a \textbf{null distribution}
from permuting stimulus labels and reindexing rows and columns together, which
says what chance looks like.

\ask{8.} The upper-triangle entries are not exchangeable --- they are generated by
$n$ stimuli, so moving one stimulus moves an entire row and column together.
Permuting entries independently destroys that structure and produces nulls far
more disordered than any real relabeling, so the null spread is too small and
every observed score looks extreme.

\ask{9.} The map is determined by the input image and the architecture, not by
anything the model learned. Any claim it supported about the trained model's
strategy is unsupported; the figure is evidence about edges.

\ask{10.} Vanishing gradients: BPTT over 300 steps multiplies 300 Jacobians, and
the product decays exponentially. Clipping bounds gradients from \emph{above} and
does nothing for gradients that are too small; the fix is gating, which routes
the cell state through an addition.

\ask{11.} A line attractor --- a direction in state space that neither decays nor
grows, with everything else contracting onto it. It supports integrating evidence
or holding a continuous value: perceptual decision-making with accumulated
evidence, or maintaining eye position or a remembered angle.

\ask{12.} Dot products of $d_k$-dimensional vectors grow like $\sqrt{d_k}$, so
without the scaling the softmax argument becomes large, the softmax saturates to
nearly one-hot, and its gradient nearly vanishes. Training then fails for large
$d_k$ while appearing fine for small $d_k$.

\ask{13.} (i) There are many heads and many layers, and their maps disagree, so
``the'' map is a selected one. (ii) Attention weight is not importance --- what a
position contributes is its value vector, and the residual stream carries
information around the attention block regardless.

\ask{14.} Exposure bias: during training the model always sees ground-truth
history, but at generation time it sees its own outputs, including its own
errors, which it was never trained to recover from. Errors compound over a
sequence.

\ask{15.} The KL term drives the encoder to match the prior, so
$\boldsymbol{\mu}\to0$, $\boldsymbol{\sigma}\to1$, and the latent code carries no
information about the input while a strong decoder reconstructs anyway.
Diagnostic: the KL term goes to (near) zero during training.

\ask{16.} Does better AI mean better models of biology? The worry is that as
benchmark performance has risen, brain-predictivity has not risen with it, and at
scale many quite different models score alike --- so the standard scores no longer
discriminate. The proposed response is to compare strategies rather than outputs.

\section*{The night before}
\addcontentsline{toc}{section}{The night before}

Eight things worth being able to state cold:

\begin{enumerate}
  \item the three splits, and what the random-labels result rules out;
  \item why sparse coding predicting V1 is a \emph{normative} explanation;
  \item selectivity from convolution, tolerance from pooling --- and the cortical
        analogue of each;
  \item the noise ceiling and the stimulus permutation, and what each is for;
  \item why a saliency map that survives weight randomization is worthless;
  \item vanishing gradients as a product, and why gating and not clipping fixes
        it;
  \item what a line attractor is, and one task that needs one;
  \item what an attention map does and does not show.
\end{enumerate}

\end{document}
